% ============================================================
% Come si usano, argomento 11: variabili continue notevoli
% ============================================================
\argomento{Le variabili continue notevoli}{Formule a p.~\pageref*{f:uniformec}}[Variabili continue notevoli]

\begin{scelta}{Come scegliere la variabile}
\renewcommand{\arraystretch}{1.25}
\begin{tabular}{@{}P{0.6\linewidth}P{0.34\linewidth}@{}}
  \toprule
  \textbf{Se la grandezza\ldots} & \textbf{la variabile è\ldots} \\
  \midrule
  cade «a caso» in un intervallo, senza punti privilegiati (fase casuale, errore di quantizzazione)
    & uniforme \vv{uniformec} \\
  è un tempo di attesa o una durata senza usura (prossimo arrivo, guasto casuale); il testo dà la media
    & esponenziale \vv{esponenziale}, con $\lambda = 1/\text{media}$ \\
  è un errore simmetrico, con una punta in 0 e code esponenziali & di Laplace \vv{laplacevar} \\
  è un rapporto o un angolo (il faro), con valori enormi ogni tanto & di Cauchy \vv{cauchyvar}: niente media \\
  è la somma di tanti piccoli contributi (rumore termico) & gaussiana (argomento 15) \\
  \bottomrule
\end{tabular}

Per ogni variabile notevole conviene ricordare densità, CDF, media e varianza: le probabilità di intervalli si fanno
con la CDF, $F_X(a_2) - F_X(a_1)$, senza integrare.
\end{scelta}

\scheda{uniformec}{
  \usoquando{Il valore cade in un intervallo e nessun punto è privilegiato: una fase casuale, l'errore di
    quantizzazione, un istante di arrivo «a caso» in una finestra.}
  \usocome{La probabilità di un sottointervallo è la sua lunghezza divisa per $b - a$.}
  \usoesempio{$X \sim \mathcal{U}(2, 6)$: $P(3 \leq X \leq 4) = \frac{1}{4}$, $E[X] = 4$, $\sigma_X^2 =
    \frac{16}{12} = \frac{4}{3}$. L'errore di quantizzazione con passo $\Delta$ è circa uniforme su
    $\left(-\frac{\Delta}{2}, \frac{\Delta}{2}\right)$: potenza $\frac{\Delta^2}{12}$.}
}

\scheda{esponenziale}{
  \usoquando{Tempi di attesa e durate senza usura: il tempo fino al prossimo pacchetto, fino a un guasto casuale.
    Se il testo dà il tempo medio, $\lambda = 1/\text{media}$.}
  \usocome{Le probabilità si leggono dalla CCDF: $P(X > x) = e^{-\lambda x}$, e $P(x_1 < X \leq x_2) =
    e^{-\lambda x_1} - e^{-\lambda x_2}$.}
  \usoattenzione{$\lambda$ non è il tempo medio ma il suo inverso: con un'attesa media di 2 minuti, $\lambda = 0{,}5$
    al minuto.}
  \usoesempio{Attesa media 2 minuti ($\lambda = 0{,}5$): $P(X > 3) = e^{-1{,}5} \simeq 0{,}2231$ e
    $P(X \leq 1) = 1 - e^{-0{,}5} \simeq 0{,}3935$.}
}

\scheda{memoria}{
  \usoquando{Si chiede la probabilità di durare «ancora $t$», sapendo che si è già durati $s$: con l'esponenziale il
    passato non conta.}
  \usoattenzione{Vale per l'esponenziale, non per qualsiasi durata: un componente che si usura «ricorda» la sua età.}
  \usoesempio{Vita media 1000 ore ($\lambda = 0{,}001$), il componente funziona da 2000 ore: la probabilità che duri
    altre 500 ore è $e^{-0{,}5} \simeq 0{,}6065$, la stessa di un componente nuovo.}
}

\scheda{laplacevar}{
  \usoquando{Errori o rumori simmetrici attorno a 0, più concentrati vicino allo zero e con code più pesanti della
    gaussiana.}
  \usocome{Per simmetria $P(X > x) = P(X < -x) = \frac{1}{2}\, e^{-\lambda x}$ per $x \geq 0$.}
  \usoesempio{$\lambda = 1$: $P(X \leq -1) = \frac{1}{2}\, e^{-1} \simeq 0{,}1839$ e
    $P(\abs{X} > 1) = e^{-1} \simeq 0{,}3679$.}
}

\scheda{cauchyvar}{
  \usoquando{Rapporti fra grandezze e angoli proiettati (il faro che illumina una parete in una direzione a caso):
    una campana simmetrica ma con valori enormi ogni tanto.}
  \usocome{Le probabilità di intervalli si fanno con la CDF e l'arcotangente.}
  \usoattenzione{Non ha media né varianza: la media campionaria non si stabilizza, e Chebyshev e la legge dei grandi
    numeri non si applicano. $a$ è solo il punto di simmetria: $F_X(a) = \frac{1}{2}$.}
  \usoesempio{$\mathcal{C}(0, 1)$: $P(-1 \leq X \leq 1) = \frac{1}{\pi}\left(\arctan 1 - \arctan(-1)\right) =
    \frac{1}{2}$ e $P(X > 1) = \frac{1}{2} - \frac{1}{\pi} \cdot \frac{\pi}{4} = \frac{1}{4}$.}
}
